\section{Introduction}
\label{sec:introduction}


The shale gas boom in the United States has received significant attention in 
economic research, particularly regarding its impact on households and local well-being. 
While the conventional wisdom suggests that resource booms lead to negative outcomes often referred to as
the Dutch Disease, characterized by a decline in other skill-intensive sectors and educational investments, this paper challenges that notion. 
We argue that the shale gas boom has had a positive impact on regional welfare in terms of income and local educational outcomes,
particularly for lower-income households.
In fact, we find that the states that experienced the shale gas boom saw a slower growth in income inequality compared to those that did not (see figure \ref{fig:gini}).
Based on a staggered difference-in-differences (DiD) framework and synthetic control method,
we suggest that the shale gas shock led to an overall increase in income and college enrollment rates among lower-income households.



In detail, we find heterogeneous effects of the shale gas discoveries in the U.S. on local income and education by household income quartiles. 
Although the shale boom did not increase household income among third and fourth quartiles, 
the income of households in the first and second income quartiles increased by approximately 4.8\% and 2.1\% respectively after three years of the shale boom.
Also, college enrollment and high school attendance in the 1Q households increased by approximately 5\%p and 18\%p in the shale boom states compared to the non-shale boom states, whereas the enrollment rates in other income quartiles did not have such salient changes.
 - It must be noted, however, that in 4Q households, the same figures slightly dropped three to four years after the shock. - 
These findings overall supplement the literature on the effects of resource booms on education
by focusing on the heterogeneous effects across different income levels and thus replacing
the insignificant or negative effects on education found in previous studies which did not address such heterogeneity.



